\documentclass{article}
\usepackage{amsthm}

\usepackage[utf8]{inputenc}     % pdfLaTeX
\usepackage[T1]{fontenc}
\usepackage[magyar]{babel}

\title{Európa}
\author{Haydu Lőrinc, Lengyel Zoltán Péter, Pintér Balázs,\\ Rovenszky Robin}
\date{Legutóbb frissítve: 2026. Április 21.}

\newtheorem{definition}{Definíció}

\begin{document}

\maketitle
\tableofcontents
\newpage

\section{Németország}

\subsection{Folyók}
\begin{itemize}
    \item Elba
    \item Rajna
    \item Duna
    \item Majna
    \item Odera
\end{itemize}

\subsection{Gazdaság}

\subsubsection{Mezőgazdaság}
\begin{itemize}
    \item GDP 1\% >
    \item Foglalkoztatottak 1\% >
    \item Ennek ellenére magas színvonal, nagy mennyiségű és jó minőségű termékek
    \begin{itemize}
        \item Korszerű/gépesített
        \item Hozzáértés
    \end{itemize}
    \item Németalföld \begin{itemize}
        \item takarmánynövények, rozs, árpa
        \item szarvasmarha
    \end{itemize}
    \item Német-középhegység \begin{itemize}
        \item Búza, kukorica
        \item sertés
    \end{itemize}
    \item Alpok \begin{itemize}
        \item Havasi pásztorkodás
    \end{itemize}
    \item Felső-Rajna vidék \begin{itemize}
        \item szőlő
    \end{itemize}
\end{itemize}

\subsubsection{Ipar}
\begin{itemize}
    \item Fő ágazatai: \begin{itemize}
        \item Vegyipar
        \item Gépipar
        \item Elektronikai ipar
    \end{itemize}
    \item Ipari szerkezetváltás okai: \begin{itemize}
        \item Kőszén szerepe visszaesett $\rightarrow$ kőszén ára lecsökkent
        \item Import nyersanyagok ára kedvező
        \item Technológiák javulása $\rightarrow$ alapanyagigények csökkenése
    \end{itemize}
    \item Ipari szerkezetváltás következménye: \begin{itemize}
        \item Modern ipari ágazatok megjelenése
        \item Nőtt a munkanélküliség
        \item Csökkentek az adóbevételek
        \item Nőtt az elvándorlók aránya
    \end{itemize}

    \item Ruhr vidék \begin{itemize}
        \item Dortmund
        \item Vegyipar: \begin{itemize}
            \item Essen
            \item Duisburg
        \end{itemize}
    \end{itemize}

    \item Megújuló energia!!!
\end{itemize}

\subsubsection{Jelentős városok}
\begin{itemize}
    \item Berlin: \begin{itemize}
        \item Gazdasági, kultúrális és politikai központ
    \end{itemize}
    
    \item München: \begin{itemize}
        \item Autógyártás (BMW)
        \item Elektronika (Siemens)
        \item Vegyipar
        \item Média
    \end{itemize}
    
    \item Stuttgart: \begin{itemize}
        \item Autógyártás (Daimler, Porsche, Mercedes)
        \item Elektronika (Bosch)
    \end{itemize}
    
    \item Frankfurt am Main: \begin{itemize}
        \item Pénzügyi központ
        \item Tőzsde
        \item Európa egyik legforgalmasabb reptere
    \end{itemize}

    \item Köln: \begin{itemize}
        \item Autógyártás (Ford)
    \end{itemize}

    \item Ludwigshafen: \begin{itemize}
        \item Vegyipar: BASF
    \end{itemize}

    \item Leverskusen: \begin{itemize}
        \item Vegyipar (Bayer)
    \end{itemize}

    \item Hamburg: \begin{itemize}
        \item Kikötőváros
        \item Kereskedelmi központ
    \end{itemize}

    \item Leipzig: \begin{itemize}
        \item Vegyipari központ
    \end{itemize}

    \item Dresden: \begin{itemize}
        \item Elektronika/szoftverfejlesztés
        \item Vegyipar
    \end{itemize}
\end{itemize}

\subsection{Közlekedési hálózat}
\begin{itemize}
    \item Közúti
    \item Vasúti: Intercity Express
    \item Vízi: \begin{itemize}
        \item Tengeri (Hamburg, Bremen, Rostock)
        \item Légi
    \end{itemize}

    \item Kombinált szállítás: környezetvédelmi okok miatt az áruk szállításának egy részét vízi vagy vasúti úton oldják meg, mert ezek környezetkímélőbb szállítási formák
\end{itemize}


\section{Ausztria és Svájc}

\subsection{Domborzat}
\begin{itemize}
    \item Aletsch-gleccser: Svájc
    \item Alpok
\end{itemize}

\subsection{Mezőgazdaság}
\begin{itemize}
    \item Növénytermesztés: völgyekben, Bécsi-medence
    \item Állattenyésztés: szarvasmarha, havasi pásztorkodás, istálló
    \item Erdőgazdálkodás: bútoripar, építőipar, papíripar \begin{itemize}
        \item Fenntartható erdőgazdálkodás
    \end{itemize}
\end{itemize}

\subsection{Energiagazdaság}
\begin{itemize}
    \item vízenergia
\end{itemize}

\subsection{Szolgáltatások}
Ausztria:
\begin{itemize}
    \item Bécs: politikai, kulturális, pénzügyi központ
    \item Burnegland
\end{itemize}
Svájc:
\begin{itemize}
    \item Bern, Genf, Zürich: bank, pénzügyek
    \item Nincs hivatalos fővárosa, de Bern aként funkcionál
\end{itemize}

\subsection{Ipar}
Ausztria:
\begin{itemize}
    \item Alapanyag gyártás: részben hazai nyersanyagok (kősó, vasérc, magnezit), részben import
    \item Gépipar \begin{itemize}
        \item Linz
        \item Steyr: mezőgazdasági gépgyártás
    \end{itemize}
    \item Fa, papír, bútoripar: Graz
\end{itemize}
Svájc:
\begin{itemize}
    \item Gépipar
    \item Vegyipar
    \item Óragyártás
    \item Lindt csoki, Toblerone
    \item Szerszámgyártás: Svájci bicska
\end{itemize}

\section{Csehország és Szlovákia}

\subsection{Csehország}
\subsubsection{Népösszetétel}
\begin{itemize}
    \item Cseh
    \item Morva
\end{itemize}

\subsubsection{Domborzat}
\begin{itemize}
    \item Érchegység
    \item Szudéták
    \item Cseh-erdő
    \item Sumava
    \item Kárpátok
    \item Cseh-Morva dombság
    \item Cseh-medence (Észak)
    \item Morva-medence (Dél)
\end{itemize}

Folyók:
\begin{itemize}
    \item Elba
    \item Moldva
\end{itemize}

\subsubsection{Mezőgazdaság}
\begin{itemize}
    \item Takarmánynövények
    \item Árpa
    \item Komló
\end{itemize}

\subsubsection{Ipar}
\begin{itemize} 
    \item Érchegység 
    
    \begin{itemize}
        \item színesfémek
        \item nemesfémek
        \item kaolin
        \item grafit
    \end{itemize}
    
    \item Szudéták, Ostrava

    \begin{itemize}
        \item kőszén
    \end{itemize}

    \item Ceruzagyártás: České Budějovice
    \item Porcelángyártás: Karlovy Vary
    \item Mezőgazdasági gépgyártás: Brno
    \item Skoda: Mlada Boleslav
\end{itemize}

\subsubsection{Szolgáltatások}
\begin{itemize}
    \item Prága pénzügyi és kulturális központ
    \item Karlovy Vary: gyógyfürdő
    \item Dél Csehország: kastélyok
    \item Sípályák
\end{itemize}

\subsection{Szlovákia}
\subsubsection{Népösszetétel}
\begin{itemize}
    \item szlovák
    \item magyar $\rightarrow$ Révkomárom, Dunaszerdahely
    \item roma
\end{itemize}

\subsubsection{Domborzat}
\begin{itemize}
    \item Kárpátok
    \item Csallóköz
\end{itemize}
Folyók:
\begin{itemize}
    \item Duna
    \item Garam
    \item Nyitra
    \item Vág
\end{itemize}

\subsubsection{Mezőgazdaság}
\begin{itemize}
    \item búza
    \item kukorica
\end{itemize}

\subsubsection{Ipar}
\begin{itemize}
    \item Kia: Zsolna
    \item Peugeot, Citroën: Nagyszombat
    \item Jaguár, Land Rover: Nyitra
    \item Volkswagen: Pozsony
    \item Kassa: vegyipar
\end{itemize}

\subsubsection{Szolgáltatások}
\begin{itemize}
    \item Pozsony kulturális és pénzügyi központ
    \item Középkori várak
    \item Sípályák
\end{itemize}

\section{Lengyelország és Ukrajna}
\subsection{Lengyelország}
\subsubsection{Nyelv, népcsoportok}

\begin{itemize}
    \item Nyugati szláv
    \item Lengyelek
\end{itemize}

\subsubsection{Domborzat}

\begin{itemize}
    \item Lengyel-alföld: jégtakaró, folyók
    \item Lengyel-középhegység
    \item Kárpátok
    \item Visztula, Odera
\end{itemize}

\subsubsection{Gazdaság}

\begin{itemize}
    \item Északon: óceáni éghajlat: \\
    takarmánynövények, rozs, árpa, burgonya, szarvasmarha
    \item Délen: nedves kontinentális égh.: \\
    +búza, kukorica
    \item Felső-Szilézia iparvidék
    \item Szudékát: kőszén
    \item Katovice: kohászat, gépgyártás
    \item Gdansk, Szcecin: vegyipar
    \item Krakkó: idegenforgalom
    \item Varsó, Lódtz: pénzügyi központok
\end{itemize}

\subsection{Ukrajna}
\subsubsection{Nyelv, népcsoportok}

\begin{itemize}
    \item Keleti szlávok
    \item Ukránok, oroszok, magyarok
\end{itemize}

\subsubsection{Domborzat}

\begin{itemize}
    \item Kelet-Európai-síkság: folyók
    \item Hátság
    \item Kárpátok
    \item Dnyeper, Donyec
\end{itemize}

\subsubsection{Gazdaság}

\begin{itemize}
    \item Száraz kontinentális (feketeföld: legkedvezőbb talajtípus): \\
    öntözéses gazdálkodás, búza, kukorica, napraforgó, cukorrépa
    \item Donyec-medence
    \item Doneck: vegyipar, kohászat, feketekőszén, vasérc
    \item Odessza: vegyipar, élelmiszeripar, vegyipar
    \item Kijev: gépgyártás, élelmiszeripar, vegyipar
\end{itemize}

\section{Románia}
\subsection{Domborzat}
\begin{definition}[Partium]
        Alföld Romániába átnyúló része. Folyók töltik fel.
\end{definition}
\begin{itemize}
    \item Moldva É-K
    \item Havasalföld D
    \item Dobrudzsa 
\end{itemize}
\begin{itemize}
    \item Keleti-Kárpátok
    \item Déli-Kárpátok
    \item Erdélyi-szigethegység
    \item Bánsági-hegyvidék
\end{itemize}
\begin{itemize}
    \item Erdélyi-medence (dombvidék)
    \item Román Alföld (folyók töltik fel)
\end{itemize}

\subsection{Folyók}
\begin{itemize}
    \item Duna (Duna-delta)
\end{itemize}

\subsection{Népesség}
\begin{itemize}
    \item Román többség
    \item Magyar kisebbség ( $\sim$1 millió)
    \item Roma kisebbség
    \item Fogyó népesség
\end{itemize}

\subsection{Gazdaság}
\subsubsection{Mezőgazdaság}
\begin{itemize}
    \item EU 3. legnagyobb mezőgazdasági termelője
    \item Népesség 1/5-ét a mezőgazdaság foglalkoztatja
    \item Román-alföld 
    \begin{itemize}
        \item Búza, kukorica (magas export), cukorrépa, napraforgó
    \end{itemize}
    \item Kárpátok
    \begin{itemize}
        \item havasi pásztorkodás - szarvasmarha, juh
    \end{itemize}
\end{itemize}

\subsubsection{Ipar}
\begin{itemize}
    \item EU-ba való belépés után rohamos fejlődésnek indult az ipar
    \item Ásványok
    \begin{itemize}
        \item Arany
        \item Ezüst
        \item Kősó
        \item Bauxit
        \item Szinesfémek
    \end{itemize}
    \item Kőolaj: Déli-Kárpátok előtere
    \item Földgáz: Erdélyi-medence
    \item Gépgyártás
    \begin{itemize}
        \item Pitesti - Dacia
        \item Brassó - Traktorok, teherautók gyártása
        \item Craiova - Ford összeszerelőüzem
    \end{itemize}
    \item Konstanca 
    \begin{itemize}
        \item Kikötő, hajógyártás, élelmiszeripar
    \end{itemize}
\end{itemize}
\subsubsection{Szolgáltatások}
\begin{itemize}
    \item Bukarest
    \item Kárpátok - Turizmus, síelés, gyógyvizek
\end{itemize}

\section{Oroszország}
\begin{itemize}
    \item A világ legnagyobb országa
    \item Európa legnépesebb fővárosa Moszkva
    \item Európához társítjuk, mert a főváros Európában helyezkedik el
    \item Itt található a világ leghosszabb vasútvonala, a Transzszibériai expressz
    \item 11 időzóna halad keresztül a területén
\end{itemize}

\subsection{Népesség} 
\begin{itemize}
    \item Közel 200 nemzetiség él Oroszország területén
    \item Legnagyobb kisebbségek az ukrán és a tatár
    \item Népsűrűség viszonylag alacsony, 8 fő/km
    \item Népesség 80\%-a Európai területen él
    \item Népesség $\frac{3}{4}$-e városokban él
\end{itemize}

\subsection{Domborzat}
\begin{itemize}
    \item Az Urál hegység ketté osztja Oroszországot \begin{itemize}
        \item Variszkuszi hegységrendszer tagja
        \item Lepusztult röghegység
    \end{itemize}
    \item Kelet-európai-síkság északi része - jég által lecsiszolt tóvidék 
    \begin{itemize}
        \item Európa legnagyobb tava - Ladoga 
    \end{itemize}
    \item Kelet-európai-síkság középső része \begin{itemize}
        \item Hátságokkal tagolt táblás vidék - üledéktakaró fedi 
    \end{itemize}
    \item Kelet-európai síkság déli része - folyók által feltöltött síkság
    \item Nyugat-szibériai alföld 
    \item Közép-szibériai fennsík
\end{itemize}

\subsection{Folyók}
\begin{itemize}
    \item Ob
    \item Jenyiszej
    \item Léna
    \item Volga
\end{itemize}

\subsection{Éghajlat}
\begin{itemize}
    \item A hideg és a mérsékelt övezet összes éghajlata megtalálható az ország területén
\end{itemize}

\subsection{Gazdaság}

\subsubsection{Mezőgazdaság}
\begin{itemize}
    \item Kelet-európai síkság 
    \begin{itemize}
        \item A mezőgazdaság fő színtere
        \item Jó minőségű, termékeny feketeföld fedi
        \item Napraforgó, búza, kukorica
    \end{itemize}
    \item Szibéria 
    \begin{itemize}
        \item Tajga éghajlat - erdőgazdálkodás
    \end{itemize}
\end{itemize}

\subsubsection{Ipar}

\begin{itemize}
    \item Kőolaj és földgáz 
    \begin{itemize}
        \item A világ legnagyobb földgázkitermelője
        \item Nagy mennyiségű kőolajat termel ki
        \item Kőolaj és földgáztermelés fő színtere: Szibéria
        \item A kőolajat és a földgézt hatalmas vezetékeken szállítják exportra
        \item A fagyott talaj miatt a kőolaj- és földgázkitermelés nehézkes, és sok környezeti kárt okoz
        \end{itemize}
    \item Iparvidékek 
    \begin{itemize}
        \item Ipar nagyrésze Európában 
        \begin{itemize}
            \item Főbb központjai: Moszkva és Szentpétervár
            \item Moszkva - gépipar, textilipar, élelmiszeripar
            \item Szentpétervár - vegyipar, elektronikai ipar, kikötőváros - kereskedés 
        \end{itemize}
        \item Urál iparvidék 
        \begin{itemize}
            \item vasérc, kálisó, nemesfémek, színesfémek 
            \item Nehézipar - kohászat
            \item Jekatyerinburg 
            \item Cseljabinszk
        \end{itemize}
    \end{itemize}
    \item A Volgára számos vízerőmű épül 
    \begin{itemize}
        \item Autógyárakat és vegyipari létesítményeket tartanak fent   
        \item Volgográd
        \item Szamara
        \item Togliatti
    \end{itemize}
\end{itemize}

\subsubsection{Szolgáltatások}
\begin{itemize}
    \item Moszkva 
    \begin{itemize}
        \item Pénzügyi, kereskedelmi és kulturális központ
    \end{itemize}
    \item Szentpétervár 
    \begin{itemize}
        \item Idegenforgalmi központ - történelmi főváros
    \end{itemize}
\end{itemize}

\section{Délszláv államok:}
\begin{itemize}
    \item Jugoszlávia $\rightarrow$ 20. század $\rightarrow$ 1991-1995 délszláv háború
    \begin{itemize}
        \item 7 önálló államra bomlott
    \end{itemize}
    \item Szlovénia, Horvátország $\rightarrow$ latin betűk
    \item Szerbia, Macedónia, Montenegró $\rightarrow$ cirill betűk, keleti kereszténység (ortodox)
    \item Bosznia-Hercegovina $\rightarrow$ iszlám vallás, latin betűk
    \item Koszovó $\rightarrow$ albán népcsoport, iszlám vallás
\end{itemize}

\subsection{Domborzat és éghajlat:}
\begin{itemize}
    \item Szlovénia $\rightarrow$ Alpok
    \item Horvátország, Bosznia-Hercegovina $\rightarrow$ Dinári-hegység $\rightarrow$ éghajlatválasztó, karsztos formák, mészkőhegység
    \item Szerb-érchegység
    \item Alföldek: Duna, Tisza, Dráva, Száva mentén 
\end{itemize}

\subsection{Gazdaság:}
\begin{itemize}
    \item Szlovénia: \begin{itemize}
        \item Legfejlettebb
        \item Legkorábban lépett be az EU-ba
        \item Háztartási gépgyártás, gyógyszer, elektronikai ipar
        \item Alpok
    \end{itemize}
    \item Horvátország: \begin{itemize}
        \item Tengerpart
        \item Plitvicei-tavak 
        \item Dalmácia: Zadar, Split, Dubrovnik
        \item Zágráb $\rightarrow$ főváros
    \end{itemize}
    \item Szerbia: \begin{itemize}
        \item Színes és nemesfémek (Szerb-érchegység) feldolgozása
        \item fejletlenebb
        \item Belgrád
    \end{itemize}
    \item Montenegró: Kotori-öböl
\end{itemize}

\end{document}